\documentclass[10pt,a4paper]{amsart}
\usepackage[margin=19mm]{geometry}
\usepackage{amsmath,amssymb,mathtools}
\usepackage[scr=rsfs,cal=euler]{mathalfa}
\usepackage{xfrac}
\usepackage[unicode,hidelinks,bookmarks=false]{hyperref}
\newcommand{\ie}{i.e.\ }
\newcommand{\cf}{cf.\ }
\newcommand{\resp}{respectively\ }
\newcommand{\Subsection}{\subsection}
\providecommand{\nobreakdash}{\nobreak\hbox{-}}
\setlength{\emergencystretch}{2em}
\multlinegap0pt
\usepackage{amssymb}
\usepackage{mathtools}
\newtheorem{proposition}{Proposition}
\newtheorem{definition}{Definition}
\title{On an Initial Value Problem Describing the Small Oscillations of a Floating Cylinder}
\author{Vicente Ocqueteau (IMB), Marius Tucsnak (IMB)}
\date{Source: arXiv:2504.20523v3 (CC BY 4.0)}
\begin{document}
\maketitle

\noindent\emph{Source and licence.} On an Initial Value Problem Describing the Small Oscillations of a Floating Cylinder. Version arXiv:2504.20523v3 (CC BY 4.0), distributed by arXiv under the Creative Commons Attribution 4.0 International licence, http://creativecommons.org/licenses/by/4.0/, as recorded in the arXiv licence metadata for this exact version and shown on the abstract page https://arxiv.org/abs/2504.20523v3. Full licence notice: \texttt{licenses/arxiv-2504.20523-CC-BY-4.0.txt}.\par\smallskip

\noindent\textit{Editorial scope.} Counted unit: the article's proposition formula\_hilbert (source0.tex lines 698--707). The article's complete original proof of it is delivered with it (source0.tex lines 709--800, one \texttt{proof} environment of 92 lines). No mathematics is written, repaired or re-derived: the statement and the proof are the article's own lines, in the article's own notation, and no retained line is substituted or re-flowed.  Retained uncounted as the source-local context the proof needs: lines 273--275; lines 380--382; lines 383--385; lines 436--439; lines 453--455; lines 548--556; lines 563--572; lines 578--584; lines 679--687.  Nothing was omitted from any retained span, so the package carries no omission record.  No retained line contains a citation, so the wrapper carries no bibliography and no bibliography entry.  No retained line contains a figure environment, an image-inclusion command or drawing code, so no image is delivered and the package declares no asset dependency.  Editorial formatting only, in addition: an AMS \texttt{amsart} preamble that restates the article's own declarations and macros used by the retained lines, namely the environments \texttt{proposition}, \texttt{definition} on separate counters, together with the standard packages those lines need.  The retained environments print in this document as Proposition 1, Definition 1 in order of appearance, whereas the article prints the same items with the numbering of its own full document.  The complete native source of the article as distributed in the arXiv e-print for this exact version is bundled unchanged as \texttt{sources/177255/source0.tex}.
\begin{equation}\label{Omega}
    \Omega=\left\{\left.\begin{bmatrix} x\\ y\end{bmatrix}\in \mathbb{R}^2\ \ \right|\ \ x^2+y^2>1,y>0\right\}.
\end{equation}

\begin{equation}\label{Diri_tot}
\Delta \tilde \psi(x,y)=0 \qquad\qquad(x\in \mathbb{R},\ y>0),
\end{equation}

\begin{equation}\label{frontiera_tot}
\tilde\psi(x,0)=\eta(x) \qquad\qquad\qquad(x\in \mathbb{R}).
\end{equation}

\begin{equation}\label{prima_parte_finita}
(\mathcal{H} f)(x)=\frac1\pi  {\rm P.V.} \int_\mathbb{R} \frac{f(\tilde x)}{x- \tilde x}\, {\rm d}\tilde x
\qquad\qquad\qquad(f\in L^2(\mathbb{R})),
\end{equation}

\begin{equation}\label{cu_Hilbert}
(\Lambda_\mathbb{H}\eta)(x)=-(\mathcal{H}\eta')(x) \qquad\qquad\qquad(\eta\in W^{1,2}(\mathbb{R})).
\end{equation}

\begin{equation}\label{equation_Dw_Laplace}
\left\{
\begin{aligned}
&\frac{\partial^2\psi}{\partial x^2}(x,y) + \frac{\partial^2\psi}{\partial y^2}(x,y) =0 && \qquad(x^2+y^2>1, y>0),\\
&\frac{\partial\psi}{\partial r}(x,y)=0 && \qquad (x\in \mathcal{E},\ y>0,\ x^2+y^2=1),\\
&\psi(x,0)=v(x) &&\qquad (|x|>1),
\end{aligned}
\right.
\end{equation}

\begin{proposition}\label{21}
For  $v\in W^{1,2}(\mathcal{E})$, the problem \eqref{equation_Dw_Laplace} has a unique solution $\psi=D_\Omega v$ in the class of functions which are continuous and bounded on $\overline\Omega$. Moreover,
  for every  $\begin{bmatrix} x\\ y\end{bmatrix}\in \Omega$  we have
\begin{equation}\label{explicit_formula_Dv}
    (D_\Omega v)(x,y)=
    \frac{1}{\pi}\int_{\mathcal{E}}\left[\frac{y}{( x-\tilde x)^2+y^2}+\frac{y}{(x\tilde x-1)^2+y^2\tilde x^2}\right]v(\tilde x)\,{\rm d}\tilde x .
\end{equation}
Finally, $D_\Omega$ is a linear bounded operator from $W^{1,2}(\mathcal{E})$ to $C_b(\overline\Omega)$, where $C_b(\overline{\Omega})$
stands for the space of continuous bounded functions on $\overline{\Omega}$, endowed with the $L^\infty$ norm.
\end{proposition}

\begin{equation}\label{numar_eta_prim}
(\eta(v))(x)=\begin{cases}
    v(x) & \qquad\qquad(x\in \mathcal{E})\\
    v\left(\frac1x\right) & \qquad\qquad(x\in (-1,1)\setminus \{0\}),\\
    0 & \qquad\qquad(x=0).
\end{cases}
\end{equation}

\begin{definition}\label{from_definition}
The Dirichlet-to-Neumann
map associated to the boundary value problem \eqref{equation_Dw_Laplace},
denoted by $\Lambda_\Omega$, is defined on $V_0(\mathcal{E})$ by
\begin{equation}\label{l_restrans}
(\Lambda_\Omega v)(x)=\left[\Lambda_\mathbb{H} (\eta(v))\right](x) \qquad\qquad\qquad(v\in V_0(\mathcal{E}), x\in \mathcal{E}),
\end{equation}
where $\eta(v)$ has been defined in \eqref{numar_eta_prim}.
\end{definition}

\begin{proposition}\label{formula_hilbert}
Let $\Omega$ be defined in  \eqref{Omega} and let $\Lambda_\Omega$ be the Dirichlet-to-Neumann map
introduced in Definition \ref{from_definition}.
Then $\Lambda_\Omega$ extends to an operator, still denoted by $\Lambda_\Omega$, in $\mathcal{L}(W^{1,2}(\mathcal{E}),L^2(\mathcal{E}))$. Moreover, we have
\begin{multline}\label{Mare_Formula}
\left(\Lambda_\Omega v\right)(x)=-(\mathcal{H}(g(v))(x) - \frac{1}{x^2} (\mathcal{H}(g(v))\left(\frac1x\right)\\
-\frac{v(1)-v(-1)}{\pi x} \qquad\qquad(v\in W^{1,2}(\mathcal{E}),\ x\in \mathcal{E}),
\end{multline}
where  $g(v)\in L^2(\mathbb{R})$ is the extension of $v'$ by zero outside $\mathcal{E}$ and $\mathcal{H}$ is the Hilbert transform defined in \eqref{prima_parte_finita}.
\end{proposition}

\begin{proof}
Let $v\in W^{1,2}(\mathcal{E})$ be such that there exists $a>0$ with $v(x)=0$ for  $|x|\geqslant a$. Then the function $\eta(v)$ defined by
\eqref{numar_eta_prim} lies in $W^{1,2}(\mathbb{R})$. For the sake of simplicity, this function will simply be denoted by $\eta$ in the remaining part of this proof.

As seen in the proof of Proposition \ref{21}
the solution $\psi$ of \eqref{equation_Dw_Laplace} is the restriction to $\overline\Omega$ of the
solution $\tilde\psi$ of   \eqref{Diri_tot}-\eqref{frontiera_tot}. We thus have
\begin{equation}\label{simetria}
\left(\Lambda_\Omega v\right)(x)=\left(\Lambda_\mathbb{H} \eta\right)(x) \qquad\qquad\qquad(x\in \mathcal{E}\ {\rm a.e.})\, .
\end{equation}
Using the fact that $\Lambda_\mathbb{H}\in \mathcal{L}(W^{1,2}(\mathbb{R}),L^2(\mathbb{R}))$ it follows that $\Lambda_\Omega v\in L^2(\mathcal{E})$.
Moreover, using \eqref{cu_Hilbert} it follows that
\begin{multline}\label{TREI_parte_finita}
\left(\Lambda_\Omega v\right)(x)= - \frac1\pi  \lim_{\varepsilon\to 0+}
\int_{|\tilde x-x|\geqslant\varepsilon} \frac{\eta'(\tilde x)}{x-\tilde x}\, {\rm d}\tilde x\\
=-\frac1\pi  \lim_{\varepsilon\to 0+} \left(
\int_{\mathcal{E}_\varepsilon(x)} \frac{\eta'(\tilde x)}{x-\tilde x}\, {\rm d}\tilde x +
\int_{\mathcal{I}_\varepsilon(x)} \frac{\eta'(\tilde x)}{x-\tilde x}\, {\rm d}\tilde x\right)
\qquad\qquad( x\in \mathcal{E}),
\end{multline}
where
\begin{equation}\label{defeeps}
 \mathcal{E}_{\varepsilon}(x)=\{\tilde x\in \mathcal{E}\ \ |\ \ |x-\tilde x|>\varepsilon\} \qquad\qquad\qquad(x\in \mathcal{E}).
 \end{equation}
\begin{equation}\label{defieps}
 \mathcal{I}_{\varepsilon}(x)=\{\tilde x\in (-1,1)\ \ |\ \ |x-\tilde x|>\varepsilon\} \qquad\qquad\qquad(x\in \mathcal{E}).
 \end{equation}
 It is easily seen that
\begin{equation}\label{defieps_noua}
 \mathcal{I}_{\varepsilon}(x)=(-1,1) \qquad\qquad\qquad(x\in \mathcal{E}, \varepsilon <|x|-1),
 \end{equation}
 so that \eqref{TREI_parte_finita} can be rewritten
\begin{equation}\label{TREI_parte_finita_bus}
\left(\Lambda_\Omega v\right)(x)= - \frac1\pi  \lim_{\varepsilon\to 0+} \left(
\int_{\mathcal{E}_\varepsilon(x)} \frac{\eta'(\tilde x)}{x-\tilde x}\, {\rm d}\tilde x +
\int_{-1}^1 \frac{\eta'(\tilde x)}{x-\tilde x}\, {\rm d}\tilde x\right)
\qquad\qquad( x\in \mathcal{E}).
\end{equation}
 We next remark that
 $$
 \lim_{\varepsilon\to 0+}
\int_{\mathcal{E}_\varepsilon(x)} \frac{\eta'(\tilde x)}{x-\tilde x}\, {\rm d}\tilde x
=\lim_{\varepsilon\to 0+}
\int_{|\tilde x-x|>\varepsilon} \frac{g(\tilde x)}{x-\tilde x}\, {\rm d}\tilde x
=\pi(\mathcal{H}g)(x) \qquad\qquad(x\in \mathcal{E}).
$$
Denote
\begin{equation}\label{defh1}
h_1(x)=\lim_{\varepsilon\to 0+}
\int_{\mathcal{E}_\varepsilon(x)} \frac{\eta'(\tilde x)}{x-\tilde x}\, {\rm d}\tilde x
\qquad\qquad(x\in \mathcal{E}).
\end{equation}
Using the fact that the Hilbert transform, as defined in \eqref{prima_parte_finita}, is an isometry of $L^2(\mathbb{R})$, we obtain that
\begin{multline}\label{est_unul_termen}
    \|h_1\|_{L^2(\mathcal{E})}=\pi\|\mathcal{H}g(v)\|_{L^2(\mathcal{E})}
    \leqslant \pi\|\mathcal{H}g(v)\|_{L^2(\mathbb{R})}\\
    =\pi \|g(v)\|_{L^2(\mathbb{R})}
    \leqslant \pi \|v\|_{W^{1,2}(\mathcal{E})}.
\end{multline}
Let
\begin{equation}\label{defh2}
h_2(x)=\int_{-1}^1 \frac{\eta'(\tilde x)}{x-\tilde x}\, {\rm d}\tilde x \qquad\qquad\qquad(x\in \mathcal{E}).
\end{equation}
Then
\begin{multline}\label{asta_e_h2}
h_2(x)=\frac{\eta(1)-\eta(-1)}{x} + \frac1x \int_{-1}^1 \frac{\tilde x \eta'(\tilde x)}{x-\tilde x}\, {\rm d}\tilde x\\
=\frac{v(1)-v(-1)}{x} - \frac{1}{x^2} \int_{-1}^1 \frac{ v'\left(\frac{1}{\tilde x}\right)}{\frac{1}{\tilde x}
-\frac{1}{x}} \frac{1}{{\tilde x}^2} \, {\rm d}\tilde x\\
=\frac{v(1)-v(-1)}{x} - \frac{1}{x^2} \int_\mathcal{E} \frac{ v'(s)}{s
-\frac{1}{x}}  \, {\rm d} s\\
=\frac{v(1)-v(-1)}{x} + \frac{1}{x^2} (\mathcal{H}g)\left(\frac1x\right) \qquad\qquad(x\in \mathcal{E}).
\end{multline}
From the last formula it clearly follows that
\begin{equation}\label{est_doi_termen}
\|h_2\|_{L^2(\mathcal{E})} \leqslant c\|v\|_{W^{1,2}(\mathcal{E})},
\end{equation}
where $c$ is a universal constant.

On the other hand, from \eqref{TREI_parte_finita_bus} we have
\begin{equation}\label{e_o_suma}
\left(\Lambda_\Omega v\right)(x)=-\frac{1}{\pi}(h_1(x) + h_2(x)) \qquad\qquad(x\in \mathcal{E}).
\end{equation}
Combining the last formula with \eqref{est_unul_termen} and \eqref{est_doi_termen} it follows that there exists a universal
constant $c>0$ such that the inequality
$$
\|\Lambda_\Omega v\|_{L^2(\mathcal{E})} \leqslant c\|v\|_{W^{1,2}(\mathcal{E})}
$$
holds for every $v\in W^{1,2}(\mathcal{E})$ which vanishes on $(-\infty,-a]\cup [a,+\infty)$ for some $a>0$.
By a density argument, it follows that the last inequality holds for every $v\in W^{1,2}(\mathcal{E})$, which implies the first conclusion of the proposition.

Finally, \eqref{Mare_Formula} follows from \eqref{defh1},\eqref{asta_e_h2} and \eqref{e_o_suma}.
\end{proof}

\end{document}
